\newpage

\section{Сопряженный Frank--Wolfe в задачах машинного обучения}
\label{sec:cfw-ml}

В этой главе модификации Frank--Wolfe с сопряженными направлениями переносятся из транспортной постановки в задачи машинного обучения.
Рассматривается гладкая выпуклая задача
\begin{equation}
    f(x)\to\min_{x\in\cX},
    \qquad
    f:\R^d\to\R,
    \label{eq:ml-constrained-problem}
\end{equation}
где $\cX$ является компактным выпуклым множеством.
Такая постановка возникает, например, когда регуляризованная задача машинного обучения переписывается как задача с ограничением на норму параметров.

Если исходная задача имеет вид
\[
    f(x)+\mu g(x)\to\min_x,
\]
то при подходящем выборе радиуса она может быть рассмотрена как
\[
    f(x)\to\min_{g(x)\leq R}.
\]
Для ограничений вида $\norm{x}_1\leq R$ линейный оракул Frank--Wolfe вычисляется проще, чем проекция на соответствующее множество.
Поэтому FW и его модификации естественно применимы к обучению моделей с ограничениями на параметры.

\subsection{Базовый FW для гладкой выпуклой задачи}

На итерации $k$ алгоритм решает линейную подзадачу
\begin{equation}
    s_k^{FW}
    =
    \argmin_{s\in\cX}
    \inner{\nabla f(x_k)}{s},
    \label{eq:ml-fw-lmo}
\end{equation}
после чего делает шаг
\[
    d_k=s_k^{FW}-x_k,
    \qquad
    x_{k+1}=x_k+\gamma_kd_k.
\]
Длина шага может задаваться теоретической последовательностью
\[
    \gamma_k=\frac{2}{k+2},
\]
или выбираться линейным поиском
\[
    \gamma_k=
    \argmin_{\gamma\in[0,1]}
    f(x_k+\gamma d_k).
\]

\begin{algorithm}
\caption{Frank--Wolfe для задачи машинного обучения}
\begin{algorithmic}
\REQUIRE компактное выпуклое множество $\cX$, начальная точка $x_0\in\cX$
\STATE $k\gets 0$
\REPEAT
\STATE $s_k^{FW}\gets\argmin_{s\in\cX}\inner{\nabla f(x_k)}{s}$
\STATE $d_k\gets s_k^{FW}-x_k$
\STATE $\gamma_k\gets\argmin_{\gamma\in[0,1]}f(x_k+\gamma d_k)$
\STATE $x_{k+1}\gets x_k+\gamma_kd_k$
\STATE $k\gets k+1$
\UNTIL{достигнут критерий остановки}
\end{algorithmic}
\end{algorithm}

\subsection{CFW без явного вычисления гессиана}

В транспортной задаче гессиан модели Бекмана диагонален, поэтому вычисление коэффициента CFW по формуле~\eqref{eq:cfw-alpha-traffic} дешево.
В задачах машинного обучения явное вычисление гессиана может быть слишком дорогим.
Поэтому используется приближение произведения гессиана на вектор через разность градиентов.

Пусть
\[
    d_{k-1}=s_{k-1}-x_{k-1},
    \qquad
    x_k=x_{k-1}+\gamma_{k-1}d_{k-1}.
\]
Тогда
\[
    s_{k-1}-x_{k-1}
    =
    \frac{x_k-x_{k-1}}{\gamma_{k-1}},
\]
и произведение гессиана на предыдущее направление приближенно выражается как
\[
    \nabla^2 f(x_k)(s_{k-1}-x_{k-1})
    \approx
    \frac{\nabla f(x_k)-\nabla f(x_{k-1})}{\gamma_{k-1}}.
\]
После сокращения $\gamma_{k-1}$ получается формула
\begin{equation}
    \alpha =
    \frac{
        (x_k-s_{k-1})^T(\nabla f(x_k)-\nabla f(x_{k-1}))
    }{
        (s_k^{FW}-s_{k-1})^T(\nabla f(x_k)-\nabla f(x_{k-1}))
    }.
    \label{eq:cfw-alpha-grad-diff}
\end{equation}
Эта формула не требует явного построения гессиана и хорошо согласуется с реализацией в библиотеках автоматического дифференцирования.

\begin{algorithm}
\caption{Conjugate Frank--Wolfe для общей гладкой задачи}
\begin{algorithmic}
\REQUIRE $x_0\in\cX$, параметр $\alpha_0\in(0,1]$
\STATE выполнить один шаг FW и получить $x_1$, $s_0$
\STATE $k\gets 1$
\REPEAT
\STATE $s_k^{FW}\gets\argmin_{s\in\cX}\inner{\nabla f(x_k)}{s}$
\IF{$\inner{\nabla f(x_k)}{s_k^{FW}}=\inner{\nabla f(x_k)}{x_k}$}
\STATE остановить алгоритм
\ENDIF
\STATE вычислить $\alpha$ по формуле~\eqref{eq:cfw-alpha-grad-diff}
\STATE $\alpha\gets\min\{1,\max\{\alpha_0,\alpha\}\}$
\STATE $s_k\gets\alpha s_k^{FW}+(1-\alpha)s_{k-1}$
\STATE $d_k\gets s_k-x_k$
\STATE $\gamma_k\gets\argmin_{\gamma\in[0,1]}f(x_k+\gamma d_k)$
\STATE $x_{k+1}\gets x_k+\gamma_kd_k$
\STATE $k\gets k+1$
\UNTIL{достигнут критерий остановки}
\end{algorithmic}
\end{algorithm}

\subsection{Сходимость CFW}

Для анализа используется стандартное предположение о выпуклости и гладкости.

\begin{assumption}
\label{ass:ml-smooth}
Множество $\cX\subset\R^d$ выпукло и компактно.
Функция $f:\R^d\to\R$ выпукла и непрерывно дифференцируема, а ее градиент липшицев на $\cX$ с константой $L$:
\[
    0
    \leq
    f(x')-f(x)-\inner{\nabla f(x)}{x'-x}
    \leq
    \frac{L}{2}\norm{x'-x}_2^2,
    \qquad
    x,x'\in\cX.
\]
\end{assumption}

Зазор Frank--Wolfe в точке $x_k$ равен
\[
    g_k=
    \inner{\nabla f(x_k)}{x_k-s_k^{FW}}.
\]
Обозначим
\[
    h_0=f(x_0)-f(x^\ast),
    \qquad
    R=\Max_{x,y\in\cX}\norm{x-y}_2,
    \qquad
    A=LR^2.
\]

\begin{lemma}
\label{lem:cfw-descent-direction}
Пусть $d_k$ является направлением шага CFW, а длина шага выбирается точным линейным поиском на отрезке $[0,1]$.
Тогда
\[
    \inner{\nabla f(x_{k+1})}{d_k}\leq 0.
\]
\end{lemma}

\begin{proof}
Докажем утверждение по индукции.
На нулевой итерации направление CFW совпадает с направлением FW:
\[
    d_0=d_0^{FW}=s_0^{FW}-x_0.
\]
Рассмотрим одномерную функцию
\[
    \varphi(\gamma)=f(x_0+\gamma d_0)
\]
на отрезке $[0,1]$.
Если $x_0$ не является решением, то по условию оптимальности для выпуклой задачи
\[
    \varphi'(0)=
    \inner{\nabla f(x_0)}{s_0^{FW}-x_0}<0.
\]
Следовательно, точный линейный поиск выбирает $\gamma_0>0$.
Так как $\varphi$ выпукла, условие оптимальности одномерной задачи
\[
    \gamma_0\in\argmin_{\gamma\in[0,1]}\varphi(\gamma)
\]
дает
\[
    \varphi'(\gamma_0)\leq 0,
    \qquad
    \inner{\nabla f(x_1)}{d_0}\leq 0.
\]
Если же $x_0$ уже является решением, алгоритм останавливается, и утверждение не требуется.

Пусть теперь для предыдущей итерации выполнено
\[
    \inner{\nabla f(x_k)}{d_{k-1}}\leq 0.
\]
Для направления CFW имеем
\[
    d_k
    =
    \alpha d_k^{FW}
    +
    (1-\alpha)(s_{k-1}-x_k)
    =
    \alpha d_k^{FW}
    +
    (1-\alpha)(1-\gamma_{k-1})d_{k-1}.
\]
Здесь $\alpha\in[\alpha_0,1]$.
По определению линейного оракула
\[
    \inner{\nabla f(x_k)}{d_k^{FW}}
    =
    \inner{\nabla f(x_k)}{s_k^{FW}-x_k}
    =
    -g_k\leq 0,
\]
а второе слагаемое не положительно по предположению индукции.
Значит,
\[
    \inner{\nabla f(x_k)}{d_k}
    \leq 0.
\]
Если $g_k>0$, то производная одномерной функции
\[
    \varphi_k(\gamma)=f(x_k+\gamma d_k)
\]
в нуле отрицательна или равна нулю только в вырожденном случае.
Точный линейный поиск по выпуклой функции снова дает
\[
    \inner{\nabla f(x_{k+1})}{d_k}
    =
    \varphi_k'(\gamma_k)
    \leq 0.
\]
Тем самым индукция завершена.
\end{proof}

\begin{theorem}
\label{thm:cfw-rate}
Пусть выполнено предположение~\ref{ass:ml-smooth}.
Тогда для CFW с параметром $\alpha_0\in(0,1]$ справедлива оценка
\[
    g_N^\ast
    =
    \Min_{1\leq k\leq N}g_k
    \leq
    \frac{\max\{2h_0,A\}}{\alpha_0\sqrt{N+1}}.
\]
\end{theorem}

\begin{proof}
Так как $\gamma_k$ выбирается точным линейным поиском, для любого $\gamma\in[0,1]$
выполнено
\[
    f(x_{k+1})
    =
    \min_{\eta\in[0,1]}f(x_k+\eta d_k)
    \leq
    f(x_k+\gamma d_k).
\]
Используя $L$-гладкость, получаем
\[
    f(x_{k+1})
    \leq
    f(x_k)
    +
    \gamma\inner{\nabla f(x_k)}{d_k}
    +
    \frac{\gamma^2L}{2}\norm{d_k}_2^2.
\]
Далее
\[
    d_k
    =
    \alpha_k d_k^{FW}
    +
    (1-\alpha_k)(1-\gamma_{k-1})d_{k-1}.
\]
По лемме~\ref{lem:cfw-descent-direction}, примененной к предыдущему шагу,
\[
    \inner{\nabla f(x_k)}{d_{k-1}}\leq 0.
\]
Следовательно,
\[
    \inner{\nabla f(x_k)}{d_k}
    \leq
    \alpha_k\inner{\nabla f(x_k)}{d_k^{FW}}
    =
    -\alpha_k g_k
    \leq
    -\alpha_0 g_k.
\]
Так как $x_k,s_k\in\cX$ и $R$ является диаметром $\cX$,
\[
    \norm{d_k}_2\leq R,
    \qquad
    A=LR^2.
\]
Следовательно,
\[
    f(x_{k+1})
    \leq
    f(x_k)
    +
    \Min_{\gamma\in[0,1]}
    \left\{
        -\alpha_0\gamma g_k
        +
        \frac{\gamma^2A}{2}
    \right\}.
\]
Минимум в правой части достигается либо во внутренней точке,
либо на правой границе отрезка.
Поэтому
\[
    f(x_{k+1})-f(x_k)
    \leq
    -\frac{\alpha_0g_k}{2}
    \min\left\{
        \frac{\alpha_0g_k}{A},
        1
    \right\}.
\]
Суммируя по $k=0,\ldots,N$ и используя $f(x^\ast)\leq f(x_{N+1})$, получаем
\[
    h_0
    \geq
    \sum_{k=0}^{N}
    \frac{\alpha_0g_k}{2}
    \min\left\{
        \frac{\alpha_0g_k}{A},
        1
    \right\}.
\]
Так как $g_N^\ast=\min_{0\leq k\leq N}g_k$, то
\[
    h_0
    \geq
    (N+1)
    \frac{\alpha_0g_N^\ast}{2}
    \min\left\{
        \frac{\alpha_0g_N^\ast}{A},
        1
    \right\}.
\]
Если $\alpha_0g_N^\ast\leq A$, то
\[
    \alpha_0g_N^\ast
    \leq
    \sqrt{\frac{2Ah_0}{N+1}}
    \leq
    \frac{\max\{2h_0,A\}}{\sqrt{N+1}}.
\]
Если $\alpha_0g_N^\ast>A$, то
\[
    \alpha_0g_N^\ast
    \leq
    \frac{2h_0}{N+1}
    \leq
    \frac{\max\{2h_0,A\}}{\sqrt{N+1}}.
\]
Объединяя два случая, получаем утверждение теоремы.
\end{proof}

Оценка отличается от стандартной оценки FW только множителем $1/\alpha_0$.
Это означает, что CFW сохраняет асимптотический порядок сходимости FW по зазору Frank--Wolfe, а практическое ускорение связано с более удачным выбором направления.

\subsection{NFW в машинном обучении}

Метод Bi-conjugate Frank--Wolfe обобщается до N-conjugate Frank--Wolfe, где новое
направление строится с учетом нескольких предыдущих направлений.
Предполагается, что предыдущие направления сопряжены между собой относительно нового
гессиана:
\[
    d_{k-m}^T\nabla^2f(x_k)d_{k-n}=0,
    \qquad
    m\neq n,\quad m,n\in\{1,\ldots,N\}.
\]
Новая точка направления записывается как выпуклая комбинация текущей точки линейного
оракула и нескольких предыдущих точек:
\[
    s_k
    =
    \alpha_0s_k^{FW}
    +
    \sum_{m=1}^{N}\alpha_ms_{k-m}.
\]
Коэффициенты выбираются так, чтобы новое направление было сопряжено с предыдущими:
\[
    d_{k-m}^T\nabla^2f(x_k)d_k=0,
    \qquad m=1,\ldots,N.
\]

Обозначим $H^k=\nabla^2f(x_k)$ и
\[
    A_{k-m}=d_{k-m}^TH^kd_k^{FW},
    \qquad
    B_{k-m}=d_{k-m}^TH^kd_{k-m}.
\]
Если на текущей итерации используется $N_{\mathrm{curr}}$ предыдущих направлений, то
сначала вычисляется
\[
    \beta_{N_{\mathrm{curr}}}
    =
    \frac{-A_{k-N_{\mathrm{curr}}}}
    {B_{k-N_{\mathrm{curr}}}(1-\gamma_{k-N_{\mathrm{curr}}})},
\]
а затем последовательно для $m=N_{\mathrm{curr}}-1,\ldots,1$:
\[
    \beta_m
    =
    \frac{-A_{k-m}}{B_{k-m}(1-\gamma_{k-m})}
    +
    \frac{\gamma_{k-m}}{1-\gamma_{k-m}}
    \sum_{n=m+1}^{N_{\mathrm{curr}}}\beta_n.
\]
После этого
\[
    \alpha_m=\beta_m\alpha_0,
    \qquad
    \alpha_0=
    \frac{1}{1+\sum_{i=1}^{N_{\mathrm{curr}}}\beta_i}.
\]

Для ML-задач точное вычисление $H^k$ обычно слишком дорого.
Поэтому, как и в CFW, используется приближение произведений гессиана через разности
градиентов.
Для предыдущего направления $d_{k-m}=s_{k-m}-x_{k-m}$ применяется оценка
\[
    \nabla f(s_{k-m})-\nabla f(x_{k-m})
    \approx
    \nabla^2f(x_{k-m})(s_{k-m}-x_{k-m}).
\]
Если гессиан за несколько итераций меняется умеренно, то отношение, входящее в формулы
для коэффициентов, можно заменить выражением
\[
    \frac{A_{k-m}}{B_{k-m}}
    \approx
    \frac{
        (\nabla f(s_{k-m})-\nabla f(x_{k-m}))^Td_k^{FW}
    }{
        (\nabla f(s_{k-m})-\nabla f(x_{k-m}))^T(s_{k-m}-x_{k-m})
    }.
\]
Так NFW сохраняет идею сопряжения с несколькими предыдущими направлениями, но не требует
явного построения матрицы Гессе.

\begin{algorithm}
\caption{N-conjugate Frank--Wolfe для задачи машинного обучения}
\label{alg:nfw-ml}
\begin{algorithmic}
\REQUIRE начальная точка $x_0\in\cX$, глубина памяти $N$, параметр $\gamma_{\max}\in[0,1]$
\STATE выполнить начальный FW-шаг и получить $x_1$, $s_0$, $N_{\mathrm{curr}}\gets 1$
\REPEAT
\STATE $s_k^{FW}\gets\argmin_{s\in\cX}\inner{\nabla f(x_k)}{s}$
\STATE вычислить коэффициенты $\alpha_0,\ldots,\alpha_{N_{\mathrm{curr}}}$
\STATE $s_k\gets\alpha_0s_k^{FW}+\sum_{i=1}^{N_{\mathrm{curr}}}\alpha_is_{k-i}$
\STATE $d_k\gets s_k-x_k$
\STATE $\gamma_k\gets\argmin_{\gamma\in[0,1]}f(x_k+\gamma d_k)$
\STATE $x_{k+1}\gets x_k+\gamma_kd_k$
\STATE обновить $N_{\mathrm{curr}}$: увеличить до $N$ или сбросить до $1$, если $\gamma_k>\gamma_{\max}$
\UNTIL{достигнут критерий остановки}
\end{algorithmic}
\end{algorithm}

\subsection{Эксперименты на логистической регрессии}

Эксперименты проводились на задачах бинарной классификации с логистической регрессией на наборах Mushrooms и MNIST.
Рассматривалась задача
\[
    \norm{w}_1\leq R,
\]
а также вариант с $l_2$-регуляризацией:
\[
    f(w)=\operatorname{logloss}(w)+\frac{\lambda}{2}\norm{w}_2^2.
\]
Интерес представляло сравнение случаев, когда решение безусловной задачи лежит внутри допустимого множества, около границы или далеко вне него.

\begin{table}[H]
\centering
\caption{$l_1$-нормы решений безусловной задачи для разных датасетов и регуляризаций.}
\begin{tabular}{lcc}
\toprule
Датасет & Без регуляризации & $l_2$-регуляризация \\
\midrule
Mushrooms & $>1000$ & $30.5$ \\
MNIST & $2.77$ & $0.15$ \\
\bottomrule
\end{tabular}
\label{tab:cfw-ml-l1-norms}
\end{table}

\noindent\textbf{Частота настоящей сопряженности.}
Так как коэффициент $\alpha$ в CFW проецируется на $[\alpha_0,1]$, важно понять,
как часто алгоритм действительно строит сопряженное направление, а не просто использует
усеченный коэффициент.
Направление считается сопряженным в этом смысле, если неусеченное значение $\alpha$ уже
лежит в отрезке $[\alpha_0,1]$.
Для численной проверки использовалась нормированная величина сопряженности
\[
    \eta_k
    =
    \frac{
        \left|d_{k-1}^T\nabla^2 f(x_k)d_k\right|
    }{
        \sqrt{
            \left(d_{k-1}^T\nabla^2 f(x_k)d_{k-1}\right)
            \left(d_k^T\nabla^2 f(x_k)d_k\right)
        }
    }.
\]
В точной арифметике для сопряженных направлений $\eta_k=0$.
Нормировка нужна, чтобы сравнивать значения на разных итерациях независимо от масштаба
направлений.
На рис.~\ref{fig:cfw-ml-conjugacy-frequency} показано значение $\eta_k$ на первых ста
итерациях CFW.
В эксперименте примерно в $80\%$ итераций новое направление действительно оказывается
сопряженным с предыдущим; при этом на начальных итерациях такие направления встречаются
реже.

\begin{figure}[H]
\centering
\includegraphics[width=0.72\textwidth]{cfw_ml/Experiment_3_fig_1.pdf}
\caption{Частота настоящей сопряженности на первых 100 итерациях CFW.}
\label{fig:cfw-ml-conjugacy-frequency}
\end{figure}

\noindent\textbf{Сравнение модификаций.}
Сравнивались FW с линейным поиском, FW с теоретическим шагом, CFW и NFW при $N=2,3,4,5$.
Во всех экспериментах для CFW использовалось $\alpha_0=0.05$.
Критерием качества был зазор Frank--Wolfe.

На рис.~\ref{fig:cfw-ml-mushrooms-12}--\ref{fig:cfw-ml-mnist-l2-34} приведены графики
сходимости для четырех групп задач: Mushrooms без регуляризации, Mushrooms с
$l_2$-регуляризацией, MNIST без регуляризации и MNIST с $l_2$-регуляризацией.
Для каждого датасета рассматриваются разные радиусы ограничения $\norm{w}_1\leq R$.

\begin{figure}[H]
\centering
\includegraphics[width=0.96\textwidth]{cfw_ml/Experiment_5_fig_12.pdf}
\caption{Mushrooms без $l_2$-регуляризации: сходимость FW, CFW и NFW по FW-зазору для первых значений радиуса $R$.}
\label{fig:cfw-ml-mushrooms-12}
\end{figure}

\begin{figure}[H]
\centering
\includegraphics[width=0.96\textwidth]{cfw_ml/Experiment_5_fig_34.pdf}
\caption{Mushrooms без $l_2$-регуляризации: сходимость FW, CFW и NFW по FW-зазору для больших значений радиуса $R$.}
\label{fig:cfw-ml-mushrooms-34}
\end{figure}

\begin{figure}[H]
\centering
\includegraphics[width=0.96\textwidth]{cfw_ml/Experiment_6_fig_12.pdf}
\caption{Mushrooms с $l_2$-регуляризацией: сходимость FW, CFW и NFW по FW-зазору для первых значений радиуса $R$.}
\label{fig:cfw-ml-mushrooms-l2-12}
\end{figure}

\begin{figure}[H]
\centering
\includegraphics[width=0.96\textwidth]{cfw_ml/Experiment_6_fig_34.pdf}
\caption{Mushrooms с $l_2$-регуляризацией: сходимость FW, CFW и NFW по FW-зазору для больших значений радиуса $R$.}
\label{fig:cfw-ml-mushrooms-l2-34}
\end{figure}

\begin{figure}[H]
\centering
\includegraphics[width=0.96\textwidth]{cfw_ml/Experiment_7_fig_12.pdf}
\caption{MNIST без $l_2$-регуляризации: сходимость FW, CFW и NFW по FW-зазору для первых значений радиуса $R$.}
\label{fig:cfw-ml-mnist-12}
\end{figure}

\begin{figure}[H]
\centering
\includegraphics[width=0.96\textwidth]{cfw_ml/Experiment_7_fig_34.pdf}
\caption{MNIST без $l_2$-регуляризации: сходимость FW, CFW и NFW по FW-зазору для больших значений радиуса $R$.}
\label{fig:cfw-ml-mnist-34}
\end{figure}

\begin{figure}[H]
\centering
\includegraphics[width=0.96\textwidth]{cfw_ml/Experiment_8_fig_12.pdf}
\caption{MNIST с $l_2$-регуляризацией: сходимость FW, CFW и NFW по FW-зазору для первых значений радиуса $R$.}
\label{fig:cfw-ml-mnist-l2-12}
\end{figure}

\begin{figure}[H]
\centering
\includegraphics[width=0.96\textwidth]{cfw_ml/Experiment_8_fig_34.pdf}
\caption{MNIST с $l_2$-регуляризацией: сходимость FW, CFW и NFW по FW-зазору для больших значений радиуса $R$.}
\label{fig:cfw-ml-mnist-l2-34}
\end{figure}

Полученные результаты показывают:
\begin{enumerate}
    \item CFW и NFW часто существенно превосходят классический FW, особенно когда решение безусловной задачи далеко от границы допустимого множества.
    \item На Mushrooms преимущество сопряженных модификаций выражено сильнее, чем на MNIST.
    \item При близости безусловного решения к границе $\cX$ различия между FW с линейным поиском и CFW могут уменьшаться.
    \item Когда безусловное решение далеко от границы допустимого множества, NFW с числом предыдущих направлений больше двух часто сходится быстрее FW и CFW.
\end{enumerate}

\noindent\textbf{Экспериментальные таблицы.}

В таблицах ниже приведены численные результаты.
В каждой строке фиксируется радиус $R$ ограничения $\norm{w}_1\leq R$.
Значения в таблице являются достигнутым FW-зазором к моменту остановки по времени;
меньшее значение означает более качественное решение.
Лучший результат в строке выделен зеленым.

\begin{table}[H]
\centering
\caption{Mushrooms без $l_2$-регуляризации. Критерий -- FW-зазор.}
\label{tab:cfw-mushrooms-full}
\scriptsize
\resizebox{\textwidth}{!}{%
\begin{tabular}{lccccccc}
\toprule
$R$ & FW ls & FW & CFW & NFW $N=2$ & NFW $N=3$ & NFW $N=4$ & NFW $N=5$ \\
\midrule
5 & $9.93\cdot10^{-4}$ & $1.54\cdot10^{-4}$ & $4.46\cdot10^{-5}$ & \textcolor{green!50!black}{$9.72\cdot10^{-10}$} & $2.11\cdot10^{-9}$ & $2.49\cdot10^{-9}$ & $5.59\cdot10^{-9}$ \\
30 & $5.11\cdot10^{-3}$ & $2.05\cdot10^{-3}$ & $1.82\cdot10^{-3}$ & $3.19\cdot10^{-4}$ & $1.40\cdot10^{-4}$ & \textcolor{green!50!black}{$7.52\cdot10^{-5}$} & $2.72\cdot10^{-4}$ \\
60 & $3.82\cdot10^{-3}$ & $2.11\cdot10^{-3}$ & $1.12\cdot10^{-3}$ & $1.14\cdot10^{-3}$ & \textcolor{green!50!black}{$3.38\cdot10^{-4}$} & $6.59\cdot10^{-4}$ & $3.83\cdot10^{-4}$ \\
150 & $9.20\cdot10^{-4}$ & $3.06\cdot10^{-2}$ & $1.34\cdot10^{-4}$ & $4.06\cdot10^{-4}$ & $1.40\cdot10^{-4}$ & \textcolor{green!50!black}{$9.27\cdot10^{-5}$} & $1.47\cdot10^{-4}$ \\
\bottomrule
\end{tabular}}
\end{table}

\begin{table}[H]
\centering
\caption{Mushrooms с $l_2$-регуляризацией. Критерий -- FW-зазор.}
\label{tab:cfw-mushrooms-l2-full}
\scriptsize
\resizebox{\textwidth}{!}{%
\begin{tabular}{lccccccc}
\toprule
$R$ & FW ls & FW & CFW & NFW $N=2$ & NFW $N=3$ & NFW $N=4$ & NFW $N=5$ \\
\midrule
5 & $1.41\cdot10^{-3}$ & $3.71\cdot10^{-4}$ & $1.56\cdot10^{-5}$ & $4.42\cdot10^{-9}$ & \textcolor{green!50!black}{$3.30\cdot10^{-9}$} & $1.86\cdot10^{-8}$ & $4.86\cdot10^{-6}$ \\
30 & $8.96\cdot10^{-3}$ & $1.33\cdot10^{-2}$ & $6.21\cdot10^{-3}$ & $7.13\cdot10^{-3}$ & $5.61\cdot10^{-3}$ & $8.16\cdot10^{-3}$ & \textcolor{green!50!black}{$5.48\cdot10^{-3}$} \\
60 & $7.11\cdot10^{-3}$ & $6.69\cdot10^{-2}$ & $2.40\cdot10^{-3}$ & $2.12\cdot10^{-3}$ & $1.57\cdot10^{-3}$ & $1.13\cdot10^{-3}$ & \textcolor{green!50!black}{$9.26\cdot10^{-4}$} \\
150 & $1.82\cdot10^{-3}$ & $5.73\cdot10^{-1}$ & \textcolor{green!50!black}{$2.06\cdot10^{-4}$} & $1.22\cdot10^{-3}$ & $1.00\cdot10^{-3}$ & $5.40\cdot10^{-4}$ & $1.34\cdot10^{-3}$ \\
\bottomrule
\end{tabular}}
\end{table}

\begin{table}[H]
\centering
\caption{MNIST без $l_2$-регуляризации. Критерий -- FW-зазор.}
\label{tab:cfw-mnist-full}
\scriptsize
\resizebox{\textwidth}{!}{%
\begin{tabular}{lccccccc}
\toprule
$R$ & FW ls & FW & CFW & NFW $N=2$ & NFW $N=3$ & NFW $N=4$ & NFW $N=5$ \\
\midrule
$0.10$ & $7.48\cdot10^{-2}$ & \textcolor{green!50!black}{$1.97\cdot10^{-2}$} & $4.00\cdot10^{-2}$ & $3.87\cdot10^{-2}$ & $2.63\cdot10^{-2}$ & $2.06\cdot10^{-2}$ & $2.98\cdot10^{-2}$ \\
$0.80$ & $5.56\cdot10^{-1}$ & $2.83$ & \textcolor{green!50!black}{$5.37\cdot10^{-1}$} & $6.45\cdot10^{-1}$ & $5.88\cdot10^{-1}$ & $6.22\cdot10^{-1}$ & $6.48\cdot10^{-1}$ \\
$1.50$ & $1.15$ & $14.17$ & \textcolor{green!50!black}{$1.04$} & $1.25$ & $1.19$ & $1.45$ & $1.48$ \\
$2.50$ & $1.93$ & $32.05$ & \textcolor{green!50!black}{$1.63$} & $2.33$ & $2.11$ & $2.44$ & $2.51$ \\
\bottomrule
\end{tabular}}
\end{table}

\begin{table}[H]
\centering
\caption{MNIST с $l_2$-регуляризацией. Критерий -- FW-зазор.}
\label{tab:cfw-mnist-l2-full}
\scriptsize
\resizebox{\textwidth}{!}{%
\begin{tabular}{lccccccc}
\toprule
$R$ & FW ls & FW & CFW & NFW $N=2$ & NFW $N=3$ & NFW $N=4$ & NFW $N=5$ \\
\midrule
$0.01$ & $3.47\cdot10^{-3}$ & $3.38\cdot10^{-4}$ & $2.63\cdot10^{-4}$ & $1.41\cdot10^{-4}$ & $1.44\cdot10^{-4}$ & \textcolor{green!50!black}{$4.47\cdot10^{-5}$} & $5.62\cdot10^{-5}$ \\
$0.10$ & $5.39\cdot10^{-2}$ & $5.75\cdot10^{-2}$ & $4.93\cdot10^{-2}$ & \textcolor{green!50!black}{$4.69\cdot10^{-2}$} & $4.92\cdot10^{-2}$ & $6.02\cdot10^{-2}$ & $5.01\cdot10^{-2}$ \\
$0.20$ & $1.63\cdot10^{-1}$ & $2.68\cdot10^{-1}$ & $1.46\cdot10^{-1}$ & $1.60\cdot10^{-1}$ & \textcolor{green!50!black}{$1.42\cdot10^{-1}$} & $1.74\cdot10^{-1}$ & $1.65\cdot10^{-1}$ \\
$0.50$ & $4.10\cdot10^{-1}$ & $2.03$ & \textcolor{green!50!black}{$3.82\cdot10^{-1}$} & $4.65\cdot10^{-1}$ & $4.28\cdot10^{-1}$ & $4.35\cdot10^{-1}$ & $4.45\cdot10^{-1}$ \\
\bottomrule
\end{tabular}}
\end{table}

\noindent\textbf{Интерпретация ML-экспериментов.}

На Mushrooms без регуляризации безусловное решение имеет очень большую $l_1$-норму.
При радиусах $R=5,30,60,150$ оно находится далеко за пределами допустимого множества.
В этом режиме траектория FW часто упирается в границу $l_1$-шара и начинает менять активные
координаты зигзагообразно.
Именно здесь сопряженные методы дают наиболее заметный выигрыш.
Особенно сильный эффект виден при $R=5$, где NFW достигает зазора порядка $10^{-9}$,
а FW остается на уровнях $10^{-4}$--$10^{-3}$.

Для Mushrooms с $l_2$-регуляризацией безусловное решение имеет норму около $30.5$.
Поэтому при $R=30$ граница близка к естественному решению задачи, а при больших радиусах
часть ограничительного эффекта ослабевает.
Тем не менее CFW и NFW сохраняют преимущество над FW.
Интересно, что при $R=150$ лучшим оказывается CFW, а не NFW.
Это показывает, что увеличение памяти не является универсально лучшим выбором:
когда активная геометрия задачи становится менее жесткой, одного предыдущего направления
может быть достаточно.

На MNIST картина слабее.
Без регуляризации оценочная $l_1$-норма безусловного решения равна примерно $2.77$,
то есть часть выбранных радиусов находится близко к решению или внутри допустимой области.
В таких условиях FW с линейным поиском и CFW/NFW становятся близкими.
Однако даже здесь CFW часто дает лучший накопленный FW-зазор, особенно при $R=0.80,1.50,2.50$.
С $l_2$-регуляризацией безусловное решение имеет норму около $0.15$.
Поэтому при $R=0.01$ ограничение жесткое, и NFW с $N=4$ оказывается лучшим.
При $R=0.50$ ограничение уже сравнительно слабое, и лучшим становится CFW.

Общий вывод экспериментов таков:
\begin{enumerate}
    \item если безусловное решение далеко от границы $\cX$, сопряженные модификации резко ускоряют FW;
    \item если безусловное решение лежит близко к границе или внутри множества, разница между методами уменьшается;
    \item для простых и разреженных задач NFW с $N>2$ часто выигрывает у CFW;
    \item для более сложных задач с плотной структурой признаков эффект сопряженности проявляется слабее;
    \item зигзагообразность не исчезает полностью, но наступает позже, чем у классического FW.
\end{enumerate}

\subsection{Выводы по CFW in ML}

CFW и NFW можно применять не только в транспортном моделировании, но и в задачах машинного обучения, если они сводятся к гладкой выпуклой оптимизации на компактном выпуклом множестве.
Для таких задач явное вычисление гессиана не является обязательным: достаточно приближения произведения гессиана на вектор через разность градиентов.
Теоретически CFW сохраняет порядок сходимости FW по зазору Frank--Wolfe.
Практически же использование предыдущих направлений позволяет уменьшить зигзагообразность траектории и ускорить обучение логистической регрессии с ограничениями на параметры.
